\documentclass[11pt,letterpaper]{article}

\usepackage[margin=1in,headheight=14pt]{geometry}
\usepackage{amsmath,amssymb,amsthm,mathtools}
\usepackage{graphicx}
\usepackage{xcolor}
\usepackage{enumitem}
\usepackage{microtype}
\usepackage{fancyhdr}
\usepackage{bm}
\usepackage[colorlinks=true,linkcolor=blue,citecolor=blue,urlcolor=blue]{hyperref}

% Personal notation and theorem environments.
%!TEX root = EM_singular_models.tex

% PDF margin etc settings
%\setlength{\textwidth}{\paperwidth}
%\addtolength{\textwidth}{-6cm}
%\setlength{\textheight}{\paperheight}
%\addtolength{\textheight}{-4cm}
%\addtolength{\textheight}{-1.1\headheight}
%\addtolength{\textheight}{-\headsep}
%\addtolength{\textheight}{-\footskip}
%\setlength{\oddsidemargin}{0.5cm}
%\setlength{\evensidemargin}{0.5cm}

%%%%%%%%%%%%%%%%%%%%%%%%%%%%%%%%

% MACROS HERE

%%%%%%%%%%%%%%%%%%%%%%%%%%%%%%%%

\newcommand{\fixedpointtheta}{\widehat\theta_\obs}

\newcommand{\cover}{N}
\newcommand{\kinit}{\nu}
\newcommand{\cdelta}{c_\threshold}
\newcommand{\dndelta}{\omega}
\newcommand{\kndelta}{\Upsilon}
\newcommand{\kndeltanew}{\Psi}
\newcommand{\radii}{\mathcal{R}}
\newcommand{\event}{\mathcal{E}}
\newcommand{\kevent}{\mathcal{D}}


% Observations, dimension etc.
\newcommand{\dims}{\ensuremath{d}}
\newcommand{\samples}{\ensuremath{n}}
\newcommand{\obs}{\samples}
\newcommand{\real}{\ensuremath{\mathbb{R}}}
\newcommand{\complex}{\ensuremath{\mathbb{C}}}
\newcommand{\realdim}{\ensuremath{\real^\dims}}

\newcommand{\smallthreshold}{\epsilon}

% some mathcal notations
\newcommand{\order}[1]{\ensuremath{\mathcal{O}\parenth{#1}}}

% Basic statistics notation
\newcommand{\xstar}{\ensuremath{x^*}}
\newcommand{\xbar}{\ensuremath{\bar{x}}}
% True parameter
\newcommand{\thetastar}{\ensuremath{{\theta^*}}}
% Estimate one
\newcommand{\thetahat}{\ensuremath{\widehat{\theta}}}
% Estimate two
\newcommand{\thetatil}{\ensuremath{\widetilde{\theta}}}
\newcommand{\thetabar}{\ensuremath{\bar{\theta}}}
\newcommand{\yvec}{\ensuremath{y}}
\newcommand{\Xmat}{\ensuremath{X}}

% Distributions
\newcommand{\distribution}{P}
\newcommand{\density}{f}
\newcommand{\gaussiandensity}{\phi}
\newcommand{\NORMAL}{\ensuremath{\mathcal{N}}}
\newcommand{\BER}{\ensuremath{\mbox{Ber}}}

% Spaces
\newcommand{\Xspace}{\ensuremath{\mathcal{X}}}
\newcommand{\Yspace}{\ensuremath{\mathcal{Y}}}
\newcommand{\Zspace}{\ensuremath{\mathcal{Z}}}


% Brackets Size
\newcommand{\brackets}[1]{\left[ #1 \right]}
\newcommand{\parenth}[1]{\left( #1 \right)}
\newcommand{\bigparenth}[1]{\big( #1 \big)}
\newcommand{\biggparenth}[1]{\bigg( #1 \bigg)}
\newcommand{\braces}[1]{\left\{ #1 \right \}}
\newcommand{\abss}[1]{\left| #1 \right |}
\newcommand{\angles}[1]{\left\langle #1 \right \rangle}
\newcommand{\tp}{^\top}

% Generic vectors and scalars
\newcommand{\myvec}{u} % for generic vector
\newcommand{\myvectwo}{v} % for generic vector
\newcommand{\mymat}{B}
\newcommand{\set}{{A}}
\newcommand{\scalar}{\alpha}
\newcommand{\tvscalar}{\epsilon}
\newcommand{\tailboundscalar}{t}
\newcommand{\term}{U}
\newcommand{\myVarone}{\alpha_X}
\newcommand{\myVartwo}{\alpha_Y}
\newcommand{\myCovmat}{\Sigma}
\newcommand{\xcord}[1]{x^{(#1)}}
\newcommand{\xcordpop}[1]{X^{(#1)}}
\newcommand{\ycordpop}[1]{Y^{(#1)}}
\newcommand{\myVarthree}{\ensuremath{U}}
\newcommand{\myVarfour}{\ensuremath{V}}
\newcommand{\myfuntwo}{\ensuremath{g}}
\newcommand{\numsteps}{t}
\newcommand{\totalnumsteps}{T}
\newcommand{\seqalpha}[1]{\alpha_{#1}}
\newcommand{\seqnu}[1]{\nu^{(#1)}}
\newcommand{\imax}{\ell_\smallthreshold}
\newcommand{\basis}{e}
\newcommand{\radem}{\varepsilon}
\newcommand{\Tone}{T\ensuremath{_{0}}}
\newcommand{\Ttwo}{T\ensuremath{_{2}}}
\newcommand{\contractPar}[1]{\ensuremath{\gamma\parenth{#1}}}

\newcommand{\simiid}{\overset{\mathrm{i.i.d.}}{\sim}}
\newcommand{\eqd}{\overset{d}{=}}


% EM updates
\newcommand{\hidden}{Z}
\newcommand{\qfun}[2]{Q(#1\vert#2)}
\newcommand{\variable}{\theta}
\newcommand{\weights}{w}
\newcommand{\myfun}{q}
\newcommand{\gradf}{\nabla \myfun}
\newcommand{\hessf}{\nabla^2 \myfun}
\newcommand{\iter}{t}
\newcommand{\mupdate}{M}
\newcommand{\mupdategeneral}{\ensuremath{M^{GN}}}




% EM contractions
\newcommand{\contractgene}{\ensuremath{\gamma^{GN}}}
\newcommand{\updateMat}{\ensuremath{\Gamma_{\theta}}}
\newcommand{\mymatTheta}{\ensuremath{B_{\theta}}}
\newcommand{\sech}{\text{sech}}

% Nhat's macros 
\newcommand{\Rspace}{\ensuremath{\mathbb{R}}}
\newcommand{\Ecal}{\ensuremath{\mathcal{E}}}
\newcommand{\Ncal}{\ensuremath{\mathcal{N}}}
\newcommand{\gradient}{\ensuremath{\nabla}}
\newcommand{\smallweight}{\ensuremath{\epsilon^{sw}}}
\newcommand{\smallnorm}{\ensuremath{\rho^{sw}}}
\newcommand{\ball}{\ensuremath{\mathbb{B}}}
\newcommand{\sphere}{\ensuremath{\mathbb{S}}}

\newcommand{\diam}{\ensuremath{\text{diam}}}
\newcommand{\upper}{\ensuremath{\bar{C}}}
\newcommand{\sample}{\ensuremath{\xi}}
\newcommand{\weight}{\ensuremath{\pi}}
\newcommand{\contracasym}{\ensuremath{\rho}}
\newcommand{\Fishermatrix}{\ensuremath{\mathcal{I}}}
\newcommand{\barcontracasym}{\ensuremath{\bar{\rho}}}


\newcommand{\TrueDensity}{\ensuremath{f}}
\newcommand{\FitDensity}{\ensuremath{f}_{\theta}}
\newcommand{\FitDensityfree}{\ensuremath{f}_{\pi, \theta}}
\newcommand{\FitDensitygen}{\ensuremath{f}_{\pi, \theta_{1}, \theta_{2}}}
\newcommand{\FitDensitymore}{\ensuremath{f}_{G}}
\newcommand{\normden}{\ensuremath{\phi}}
\newcommand{\mean}{\ensuremath{\theta}}
\newcommand{\sd}{\ensuremath{\sigma}}
\newcommand{\parFOS}{\ensuremath{\gamma}}
\newcommand{\funQ}{\ensuremath{Q}}
\newcommand{\updateM}{\ensuremath{M}}
\newcommand{\weightFun}{\ensuremath{w}}
\newcommand{\mydefn}{\ensuremath{:=}}
\newcommand{\poincare}{\ensuremath{\tilde{C}}}
\newcommand{\paraspace}{\ensuremath{\Theta}}
\newcommand{\loglihood}{\ensuremath{F}}
\newcommand{\prior}{\ensuremath{\pi}}
\newcommand{\posterior}{\ensuremath{\Pi}}
\newcommand{\boundcons}{\ensuremath{B}}
\newcommand{\boundconstot}{\ensuremath{H}}
\newcommand{\noise}{\ensuremath{\varepsilon}}
\newcommand{\loglihoodgaus}{\ensuremath{F^{G}}}
\newcommand{\loglihoodind}{\ensuremath{F^{I}}}
\newcommand{\loglihoodlogit}{\ensuremath{F^{R}}}
\newcommand{\noiseone}{\ensuremath{\varepsilon_{1}}}
\newcommand{\noisetwo}{\ensuremath{\varepsilon_{2}}}

%%%SDE_Bayesian_inf
\newcommand{\weakcon}{\ensuremath{\psi}}
\newcommand{\perturb}{\ensuremath{\zeta}}
\newcommand{\inver}{\ensuremath{\xi}}

% Universal constants
\newcommand{\unicon}{c}
\newcommand{\uniconnew}{c'}
\newcommand{\uniconone}{c_1}
\newcommand{\unicontwo}{c_2}

% some parameters that you may define
\newcommand{\scparam}{m} % strong  convexity parameter
\newcommand{\smoothness}{L} % smoothness parameter
\newcommand{\step}{h}


%%%%%%%%% Distributions and Random variables %%%%%%%%%%%
\newcommand{\rvg}{\xi} % to denote the random variable g
\newcommand{\threshold}{\delta}


%%%%%%%%% Basic Terms like defn, etal, tmix, polylog %%%%%%%%%%%
\newcommand{\defn}{:=}
\newcommand{\rdefn}{=:}
\newcommand{\etal}{{et al.}}
\newcommand{\polylog}{\text{poly-log}}

% Some vector/matrix norms
\newcommand{\matnorm}[3]{|\!|\!| #1 | \! | \!|_{{#2}, {#3}}}
\newcommand{\matsnorm}[2]{|\!|\!| #1 | \! | \!|_{{#2}}}
\newcommand{\vecnorm}[2]{\left\| #1\right\|_{#2}}
\newcommand{\enorm}[1]{\vecnorm{#1}{2}} % euclidean norm
\newcommand{\nucnorm}[1]{\ensuremath{\matsnorm{#1}{\footnotesize{\mbox{nuc}}}}}
\newcommand{\fronorm}[1]{\ensuremath{\matsnorm{#1}{\footnotesize{\mbox{fro}}}}}
\newcommand{\opnorm}[1]{\ensuremath{\matsnorm{#1}{\tiny{\mbox{op}}}}}
%%% EM paper
\newcommand{\constrainnorm}[1]{\vecnorm{#1}{[k]}} % euclidean norm

\newcommand{\bigmatsnorm}[2]{{\Big|\!\Big|\!\Big| #1
  \Big|\!\Big|\!\Big|}_{#2}}

% Inner product
\newcommand{\inprod}[2]{\ensuremath{\langle #1 , \, #2 \rangle}}

% Kullback-Leibler
\newcommand{\kull}[2]{\ensuremath{D_{\text{KL}}(#1\; \| \; #2)}}

% Probability
\newcommand{\Exs}{\ensuremath{{\mathbb{E}}}}
\newcommand{\Prob}{\ensuremath{{\mathbb{P}}}}

%Eigenvector / eigenvalue related notation
\newcommand{\eig}[1]{\ensuremath{\lambda_{#1}}}
\newcommand{\eigmax}{\ensuremath{\eig{\max}}}
\newcommand{\eigmin}{\ensuremath{\eig{\min}}}

% \DeclareMathOperator{\det}{det}
\DeclareMathOperator{\modd}{mod}
\DeclareMathOperator{\diag}{diag}
\DeclareMathOperator{\var}{var}
\DeclareMathOperator{\cov}{cov}
\DeclareMathOperator{\trace}{trace}
\DeclareMathOperator{\abs}{abs}
\DeclareMathOperator{\floor}{floor}
\DeclareMathOperator{\vol}{vol}
\DeclareMathOperator{\child}{child}
\DeclareMathOperator{\parent}{parent}
\DeclareMathOperator{\sign}{sign}
\DeclareMathOperator{\rank}{rank}
\DeclareMathOperator{\toeplitz}{toeplitz}




\newtheoremstyle{named}{}{}{\itshape}{}{\bfseries}{.}{.5em}{\thmnote{#3's }#1}
\theoremstyle{named}
\newtheorem*{namedtheorem}{Lemma}

%%%%%%%
\theoremstyle{plain}

% {Theorem, Proposition, Lemma, Corollary} numbered sequentially
% throughout the paper
\newtheorem{theorem}{Theorem}
\newtheorem{proposition}{Proposition}
\newtheorem{lemma}{Lemma}
\newtheorem{claim}{Claim}
\newtheorem{corollary}{Corollary}
\newtheorem{definition}{Definition}
\newtheorem{conjecture}{Conjecture}
\newtheorem{remark}{Remark}
%%%%%%%%%%%%%%%%%%%%%%%%%%%%%%%%%%%%%%%%%%%%%%%%%%%%%%%%%%%%%%%%%%%%%%%
% WIDEBAR COMMAND
\newlength{\widebarargwidth}
\newlength{\widebarargheight}
\newlength{\widebarargdepth}
\DeclareRobustCommand{\widebar}[1]{%
  \settowidth{\widebarargwidth}{\ensuremath{#1}}%
  \settoheight{\widebarargheight}{\ensuremath{#1}}%
  \settodepth{\widebarargdepth}{\ensuremath{#1}}%
  \addtolength{\widebarargwidth}{-0.3\widebarargheight}%
  \addtolength{\widebarargwidth}{-0.3\widebarargdepth}%
  \makebox[0pt][l]{\hspace{0.3\widebarargheight}%
    \hspace{0.3\widebarargdepth}%
    \addtolength{\widebarargheight}{0.3ex}%
    \rule[\widebarargheight]{0.95\widebarargwidth}{0.1ex}}%
  {#1}}



%%% New version of \caption puts things in smaller type, single-spaced
%%% and indents them to set them off more from the text.
\makeatletter
\long\def\@makecaption#1#2{
        \vskip 0.8ex
        \setbox\@tempboxa\hbox{\small {\bf #1:} #2}
        \parindent 1.5em  %% How can we use the global value of this???
        \dimen0=\hsize
        \advance\dimen0 by -3em
        \ifdim \wd\@tempboxa >\dimen0
                \hbox to \hsize{
                        \parindent 0em
                        \hfil
                        \parbox{\dimen0}{\def\baselinestretch{0.96}\small
                                {\bf #1.} #2
                                %%\unhbox\@tempboxa
                                }
                        \hfil}
        \else \hbox to \hsize{\hfil \box\@tempboxa \hfil}
        \fi
        }
\makeatother



%% COMMENTING commands

\long\def\comment#1{}
\definecolor{battleshipgrey}{rgb}{0.52, 0.52, 0.51}
\definecolor{darkgray}{rgb}{0.66, 0.66, 0.66}
\definecolor{darkgreen}{rgb}{0.0, 0.2, 0.13}
\definecolor{darkspringgreen}{rgb}{0.09, 0.45, 0.27}
\definecolor{dukeblue}{rgb}{0.0, 0.0, 0.61}
\definecolor{olivedrab7}{rgb}{0.24, 0.2, 0.12}
\definecolor{darkblue}{rgb}{0.0, 0.0, 0.55}
\definecolor{darkscarlet}{rgb}{0.34, 0.01, 0.1}
\definecolor{candyapplered}{rgb}{1.0, 0.03, 0.0}
\definecolor{ao(english)}{rgb}{0.0, 0.5, 0.0}
\definecolor{applegreen}{rgb}{0.55, 0.71, 0.0}

\newcommand{\red}[1]{\textcolor{red}{#1}}
\newcommand{\mjwcomment}[1]{{\bf{{\red{{MJW --- #1}}}}}}
\newcommand{\mwlcomment}[1]{{\bf{{\color{orange}{{Wenlong --- #1}}}}}}

\newcommand{\blue}[1]{\textcolor{blue}{#1}}
\newcommand{\green}[1]{\textcolor{green}{#1}}
\newcommand{\highlight}[1]{\textcolor{applegreen}{#1}}

% comment lines
\newcommand{\genecomment}[1]{{\bf{{\blue{{TEAM --- #1}}}}}}
\newcommand{\rrdcomment}[1]{{\bf{{\blue{{RRD --- #1}}}}}}
\newcommand{\nhcomment}[1]{{\bf{{\blue{{NH --- #1}}}}}}
\newcommand{\kkcomment}[1]{{\bf{{\darkgreen{{KK --- #1}}}}}}


\newcommand{\widgraph}[2]{\includegraphics[keepaspectratio,width=#1]{#2}}


\newcommand{\coursenumber}{STA355F}
\newcommand{\coursetitle}{Theory of Statistical Practice}
\newcommand{\courseterm}{Fall 2026}
\newcommand{\lecturenumber}{2}
\newcommand{\lecturedate}{September 18, 2026}
\newcommand{\lecturer}{Wenlong Mou}

\pagestyle{fancy}
\fancyhf{}
\fancyhead[L]{\coursenumber: \coursetitle}
\fancyhead[R]{Lecture \lecturenumber}
\fancyfoot[C]{\lecturenumber--\arabic{page}}
\renewcommand{\headrulewidth}{0.4pt}
\renewcommand{\footrulewidth}{0pt}
\newtheorem{example}{Example}
\setcounter{section}{0}
\renewcommand{\thepage}{\lecturenumber--\arabic{page}}

\newcommand{\makelecturenotesheader}{%
  \noindent
  \textbf{\coursenumber: \coursetitle}%
  \hfill
  \textbf{\courseterm}

  \begin{center}
    {\Large\bfseries Lecture \lecturenumber: \lecturedate}
  \end{center}

  \noindent
  \textbf{Lecturer:} \lecturer
  \par\medskip
  \hrule
  \bigskip
}

\begin{document}

\makelecturenotesheader

\section{Additional probability review}
\subsection{Probabilistic convergence and arithmetic operations}
\begin{theorem}
  Suppose that $(X_n)_{n \geq 1}$ and $(Y_n)_{n \geq 1}$ are sequences of random variables, and $X$, $Y$ are random variables.
  \begin{enumerate}
    \item If $X_n \xrightarrow{p} X$ and $Y_n \xrightarrow{p} Y$, then $X_n + Y_n \xrightarrow{p} X + Y$ and $X_n Y_n \xrightarrow{p} XY$.
    \item If $X_n \xrightarrow{\mathbb{L}^p} X$ and $Y_n \xrightarrow{\mathbb{L}^p} Y$, then $X_n + Y_n \xrightarrow{\mathbb{L}^p} X + Y$, but this does not necessarily hold for the product $X_n Y_n$.
    \item (Slutsky's theorem) If $X_n \xrightarrow{d} X$ and $Y_n \xrightarrow{p} c$ for some constant $c$, then $X_n + Y_n \xrightarrow{d} X + c$, $X_n Y_n \xrightarrow{d} cX$, and $X_n / Y_n \xrightarrow{d} X / c$ if $c \neq 0$.
  \end{enumerate}
\end{theorem}
Note: Slutsky's theorem does not hold true if $Y_n$ converges in distribution to a non-degenerate random variable $Y$ instead of a constant. On the other hand, it does not require any dependence structure between $X_n$ and $Y_n$.

Intuition about Slutsky: $X_n \xrightarrow{d} X$ only tells you about the marginal distribution of $X_n$. It does not tell you anything about the joint distribution of $(X_n, Y_n)$. But if $Y_n \xrightarrow{p} c$, you are in a degenerate case, and the joint distribution of $(X_n, Y_n)$ does not matter.


\subsection{\texorpdfstring{$o_p ()$ and $O_p ()$ notation}{op and Op notation}}
Notations about limit of sequences: let $(a_n)_{n \geq 1}$ and $(b_n)_{n \geq 1}$ be sequences of real numbers. $b_n > 0$ for all $n$.
\begin{itemize}
  \item We say that $a_n = o(b_n)$ if $\lim_{n \to \infty} a_n / b_n = 0$.
  \item We say that $a_n = O(b_n)$ if there exists a constant $C > 0$ such that $|a_n| \leq C b_n$ for all sufficiently large $n$.
\end{itemize}
e.g. $a_n = o(1)$ means $a_n \to 0$, and $a_n = O(1)$ means $a_n$ is bounded.

\noindent e.g. when $a_n = 2 n^2 + n +1 $ and $b_n = n^2$, we have $a_n = O (b_n)$.

\noindent e.g. when $a_n = 1/n$ and $b_n = \log n / n$, we have $a_n = o(b_n)$.

\paragraph{Probablistic version:} Let $(X_n, Y_n)_{n \geq 1}$ be sequences of random variables, and $Y_n > 0$ with probability 1 for all $n$.
\begin{itemize}
  \item We say that $X_n = o_p(Y_n)$ if $X_n / Y_n \xrightarrow{p} 0$.
  \item We say that $X_n = O_p (Y_n)$ if for any $\varepsilon > 0$, there exists a constant $C_\varepsilon > 0$ such that $\Prob(|X_n| \leq C_\varepsilon Y_n) \geq 1 - \varepsilon$ for all $n$. ($X_n/ Y_n$ is asymptotically tight).
\end{itemize}
Remarks:

$o_p (1)$ means converges in probability to 0, and $O_p (1)$ means bounded in probability (probabilistic mass is not escaping to infinity as $n \to \infty$).

Facts:
\begin{itemize}
  \item $O_P(1) + O_P(1) = O_P(1)$.
  \item $o_P(1) + o_P(1) = o_P(1)$.
  \item $O_P(1) \times o_P (1) = o_P(1)$.
  \item $O_P(1) \times O_P(1) = O_P (1)$.
\end{itemize}

\subsection{Delta method}
Motivation: the estimators in practice are not always sample average, but some nonlinear functions related to sample average. We want to understand the asymptotic distribution of these estimators.
\begin{theorem}[Delta method]
  Suppose that $(X_n)_{n \geq 1}$ is a sequence of random variables such that $\sqrt{n}(X_n - \mu) \xrightarrow{d} N(0, \sigma^2)$ for some constant $\mu$ and $\sigma^2 > 0$. ($X_n$ themselves do not need to be sample averages.) Let $g: \real \to \real$ be a continuously differentiable function at $\mu$ with $g'(\mu) \neq 0$. Then we have
  \begin{align*}
    \sqrt{n} \big(g(X_n) - g(\mu)\big) \xrightarrow{d} N(0, \sigma^2 (g'(\mu))^2).
  \end{align*}
\end{theorem}
Proof idea: we use Taylor expansion of $g$ at $\mu$:
\begin{align*}
  g(X_n) = g(\mu) + g'(\mu)(X_n - \mu) + o_p (|X_n - \mu|).
\end{align*}
By assumption, we have
\begin{align*}
  \sqrt{n} g'(\mu)(X_n - \mu) \xrightarrow{d} N(0, \sigma^2 (g'(\mu))^2).
\end{align*}
So we just need to deal with the residual. Since $\sqrt{n} (X_n - \mu) = O_p(1)$, we have $o_p(|X_n - \mu|) = o_p(1/\sqrt{n})$. Therefore, $\sqrt{n} o_p(|X_n - \mu|) = o_p(1) \xrightarrow{p} 0$. By Slutsky's theorem, we conclude that
\begin{align*}
  \sqrt{n} \big(g(X_n) - g(\mu)\big) = \sqrt{n} g'(\mu)(X_n - \mu) + \sqrt{n} o_p(|X_n - \mu|) \xrightarrow{d} N(0, \sigma^2 (g'(\mu))^2).
\end{align*}
Note: each steps in above proof can be made more rigorous by expanding the definitions of $o_p$ and $O_p$.

\bigskip
\noindent \textbf{Remarks:}

\begin{itemize}
  \item Continuous differentiability is necessary. e.g., when $g(x) = |x|$, and $\mu = 0$. Then we have
  \begin{align*}
    \sqrt{n} (g(X_n) - g(\mu)) = \sqrt{n} |X_n| = |\sqrt{n} X_n| \xrightarrow{d} |Y| \text{ where } Y \sim N(0, \sigma^2).
  \end{align*}
  \item If $g' (\mu) = 0$, then the middle term in the Taylor expansion vanishes, and we have $g(X_n) - g(\mu) = o_p(|X_n - \mu|) = o_p(1/\sqrt{n})$. Technically, the Delta method result is still valid, as $\sqrt{n} (g(X_n) - g(\mu)) = o_p(1) \xrightarrow{d} 0$, which is a degenerate normal distribution. But this is not useful for characterizing the shape of the distribution. In this case, we need to use higher-order Taylor expansion to get a non-degenerate limiting distribution.
\end{itemize}
\begin{theorem}[second-order Delta method]
  Suppose that $(X_n)_{n \geq 1}$ is a sequence of random variables such that $\sqrt{n}(X_n - \mu) \xrightarrow{d} N(0, \sigma^2)$ for some constant $\mu$ and $\sigma^2 > 0$. Let $g: \real \to \real$ be a twice continuously differentiable function at $\mu$ with $g'(\mu) = 0$ and $g''(\mu) \neq 0$. Then we have
  \begin{align*}
   \frac{2 n (g (X_n) - g (\mu))}{g'' (\mu) \sigma^2} \xrightarrow{d} \chi^2_1.
  \end{align*}
  Here $\chi^2_1$ is the probability distribution of $Y^2$ where $Y \sim N(0,1)$.
\end{theorem}
Proof is similar: use second-order Taylor expansion of $g$ at $\mu$:
\begin{align*}
  g(X_n) = g(\mu) + \frac{1}{2} g''(\mu) (X_n - \mu)^2 + o_p(|X_n - \mu|^2).
\end{align*}
We just use the fact
\begin{align*}
  n (X_n - \mu)^2 = (\sqrt{n} (X_n - \mu))^2 \xrightarrow{d} \sigma^2 Y^2 \text{ where } Y \sim N(0, 1).
\end{align*}
In principle, we can extend this to higher-order Delta method. But in practice, we usually only need first-order or second-order Delta method.

\paragraph{Multivariate extension:} Suppose that $(X_n)_{n \geq 1}$ is a sequence of random vectors in $\real^d$ such that $\sqrt{n}(X_n - \mu) \xrightarrow{d} N(0, \Sigma)$ for some constant vector $\mu \in \real^d$ and positive definite covariance matrix $\Sigma$. Let $g: \real^d \to \real$ be a continuously differentiable function at $\mu$ with gradient $\nabla g(\mu) \neq 0$. Then we have
\begin{align*}
  \sqrt{n} (g(X_n) - g(\mu)) \xrightarrow{d} N(0, \nabla g(\mu)^\top \Sigma \nabla g(\mu)).
\end{align*}
Remark: we can also extend this to vector-valued functions $g: \real^d \to \real^k$ with Jacobian matrix $\nabla g(\mu) \in \real^{k \times d}$, and the limiting distribution is multivariate normal with covariance matrix $\nabla g(\mu) \Sigma \nabla g(\mu)^\top \in \real^{k \times k}$.

Remark: we can also extend the second-order Delta method to multivariate case, using Taylor expansion with Hessian matrix $\nabla^2 g(\mu) \in \real^{d \times d}$. In such case, the limiting random variable is a quadratic form of a multivariate normal random vector, which is a generalized version of chi-squared distribution.


\section{Statistical models and statistical inference}
A ``statistical model'' is a class of probability distributions.
\begin{itemize}
  \item In probability theory, we work with a fixed and known probability distributions, and study its properties.
  \item In statistics, we work with a class of probability distributions, and we do not know which distribution in the class is the true one. We want to learn about the true distribution from data.
\end{itemize}
Roughly speaking, statistical models can be divided into two categories: parametric models and nonparametric models.
\begin{itemize}
  \item Parametric models: the class of distributions is indexed by a finite-dimensional parameter $\theta \in \Theta \subseteq \real^d$. For example, the class of normal distributions $\{N(\mu, \sigma^2): \mu \in \real, \sigma^2 > 0\}$ is a parametric model with parameter $\theta = (\mu, \sigma^2) \in \real \times (0, \infty)$.
  \item A nonparametric model is also ``parametrized'', but the parameter is infinite-dimensional. For example, the class of all distributions on $\real$ is a nonparametric model, where the parameter is the cumulative distribution function $F: \real \to [0,1]$ that is non-decreasing and right-continuous, which is an infinite-dimensional object.
\end{itemize}
More examples of nonparametric models:
\begin{itemize}
\item The set of all probability density functions $p$ on $[0, 1]$ such that $|p' (x)| \leq 1$ for all $x \in [0, 1]$. (we will revisit this at the end of this semester).
\item The set of all probability distribution $\Prob$ on $\real$ such that $\Exs_\Prob[|X|^2] < + \infty$.
\end{itemize}

\subsection{Tasks in statistical inference}
Overall, the goal of statistics is to learn about ``which distribution in the class is the true one'' from data. There are three main tasks in statistical inference:
\begin{enumerate}
  \item Point estimation.
  \item Confidence intervals.
  \item Hypothesis testing.
\end{enumerate}

\subsubsection{Point estimation}
In general, given a statistical model $\{\Prob_\theta: \theta \in \Theta\}$, and a functional $T: \Theta \to \real^k$ that we want to estimate, we want to construct an estimator $\widehat{T}$ based on samples such that $\vecnorm{\widehat{T} - T(\theta^*)}{2}$ is small, where $\theta^*$ is the true parameter.

Criteria for evaluating estimators: mean-squared error.
\begin{align*}
  \text{MSE}_\theta (\widehat{T}) = \Exs_\theta \Big[\vecnorm{\widehat{T} - T(\theta)}{2}^2 \Big].
\end{align*}
Here the notation $\vecnorm{x}{2}$ means the Euclidean norm of a vector $x \in \real^k$.

In general, depending on the aspects of estimators that we want to evaluate, we may also use other norms and other loss functions.

\paragraph{Bias-variance decomposition} For MSE, we have
\begin{align*}
  MSE_\theta (\widehat{T}) = \Exs_\theta \Big[\vecnorm{\widehat{T} - T(\theta)}{2}^2 \Big] = \Vert \underbrace{\Exs_\theta[\widehat{T}] - T(\theta)}_{\mathrm{bias}}\Vert_{2}^2 + \underbrace{\Exs_\theta \Big[\vecnorm{\widehat{T} - \Exs_\theta[\widehat{T}]}{2}^2\Big]}_{\mathrm{variance}}.
\end{align*}
We say that an estimator $\widehat{T}$ is unbiased if $\Exs_\theta[\widehat{T}] = T(\theta)$ for all $\theta \in \Theta$. In this case, the MSE is the variance of the estimator.

In general, we want to make the optimal trade-off between bias and variance to minimize MSE. This is also related to model selection in machine learning.

e.g. $X_1, X_2, \cdots, X_n \simiid \Prob$, and $\Exs_\Prob [X^2] < + \infty$. Suppose that we want to estimate the parameters
\begin{align*}
\mu (\Prob) = \Exs_\Prob[X], \quad \sigma^2 (\Prob) = \var_\Prob(X).
\end{align*}
(In the notation above, $\Prob$ corresponds to $\theta$, and $(\mu(\Prob), \sigma^2(\Prob))$ corresponds to $T(\theta)$.)

Canonical choices:
\begin{align*}
  \widehat{\mu}_n = \frac{1}{n} \sum_{i=1}^n X_i, \quad \widehat{\sigma}_n^2 = \frac{1}{n - 1} \sum_{i=1}^n (X_i - \widehat{\mu}_n)^2.
\end{align*}
Straightforward calculations show that $\widehat{\mu}_n$ and $\widehat{\sigma}_n^2$ are unbiased estimators of $\mu(\Prob)$ and $\sigma^2(\Prob)$, respectively.

Remark: we can also normalize $\widehat{\sigma}_n^2$ by $n$ instead of $n-1$. In this case, the estimator is biased, but it may have smaller variance. In general, we want to make the optimal trade-off between bias and variance to minimize MSE. Asymptotically, normalizing by $n$ or $n-1$ does not matter, as the difference is negligible when $n$ is large.

Important properties of estimators:
\begin{definition}
  Given $X_1, X_2, \ldots, X_n \simiid \Prob$, an estimator $\widehat{T}_n$ is said to be \emph{consistent} if $\widehat{T}_n \xrightarrow{p} T(\theta^*)$ as $n \to \infty$, where $\theta^*$ is the true parameter.
\end{definition}

\begin{definition}
  An estimator $\widehat{T}_n$ is said to be \emph{asymptotically normal} if there exists an increasing sequence of positive numbers $(a_n)_{n \geq 1}$ such that $a_n \to \infty$ as $n \to \infty$, and
  \begin{align*}
    a_n (\widehat{T}_n - T(\theta^*)) \xrightarrow{d} N(0, \Sigma),
  \end{align*}
  where $\Sigma$ is a positive definite covariance matrix.
\end{definition}
Consistency is a basic requirement, and asymptotic normality is a stronger property that allows us to quantify uncertainty of the estimator.

\subsubsection{Confidence sets}
\begin{definition}
 For fixed $\alpha \in (0,1)$, a $(1 - \alpha)$-confidence set for $T(\theta)$ is a random set $C_n$ such that
  \begin{align*}
    \Prob_\theta(C_n \ni T(\theta)) \geq 1 - \alpha, \quad \forall \theta \in \Theta.
  \end{align*}
\end{definition}
Note: $\theta$ here is NOT a random variable. The probability is taken with respect to the randomness of the data, which is used to construct the confidence set $C_n$. A random subset of $\real^k$ that contains the deterministic quantity $T(\theta)$ with probability at least $1 - \alpha$, for every possible $\theta \in \Theta$.

\begin{definition}
 Given $X_1, X_2, \ldots, X_n \simiid \Prob$, a $(1 - \alpha)$ asymptotic confidence set for $T(\theta)$ is a random set $C_n$ such that
  \begin{align*}
    \lim\inf_{n \to \infty} \Prob_\theta(C_n \ni T(\theta)) \geq 1 - \alpha, \quad \forall \theta \in \Theta.
  \end{align*}
\end{definition}

Example: let $X_1, X_2, \cdots, X_n \simiid \mathrm{Ber} (\theta)$, for unknown $\theta \in (0,1)$.

Different choices of (asymptotic) confidence intervals for $\theta$:
\begin{itemize}
  \item CI via asymptotic normality. From the corollary of CLT we have learned in the previous lecture, we have
  \begin{align*}
    \frac{\sqrt{n} (\bar{X}_n - \theta)}{\widehat{se}_n} \xrightarrow{d} \mathcal{N} (0, 1),
  \end{align*}
  where $\widehat{se}_n^2 = \frac{1}{n} \sum_{i=1}^n (X_i - \bar{X}_n)^2$ is the sample standard error.

  We can then construct an asymptotic $(1 - \alpha)$ confidence interval for $\theta$ as
  \begin{align*}
    C_n = \Big[\bar{X}_n - z_{1 - \alpha/2} \cdot \widehat{se}_n / \sqrt{n} , \bar{X}_n + z_{1 - \alpha/2} \cdot \widehat{se}_n / \sqrt{n}\Big],
  \end{align*}
  where $z_{1 - \alpha/2}$ is the $(1 - \alpha/2)$-quantile of the standard normal distribution. This is asymptotically valid (and asymptoically exact):
  \begin{align*}
    \lim_{n \to \infty} \Prob_\theta(C_n \ni \theta) = 1 - \alpha, \quad \forall \theta \in (0,1).
  \end{align*}
  On the other hand, for finite $n$, this confidence interval may not have the desired coverage probability $1 - \alpha$.
  \item CI via probability deviation inequalities. e.g., we can use Chebyshev's inequality
  \begin{align*}
    \Exs_\theta[(\bar{X}_n - \theta)^2] = \frac{\theta (1 - \theta)}{n} \leq \frac{1}{4 n}.
  \end{align*}
  By Chebyshev's inequality, we have
  \begin{align*}
    \Prob_\theta(|\bar{X}_n - \theta| \geq \varepsilon) \leq \frac{1}{4 n \varepsilon^2}.
  \end{align*}
  For this to be less than $\alpha$, we can choose $\varepsilon = \sqrt{1 / (4 n \alpha)}$. This gives a finite-sample valid confidence interval for $\theta$:
  \begin{align*}
    C_n = \Big[\bar{X}_n - \sqrt{\frac{1}{4 n \alpha}}, \bar{X}_n + \sqrt{\frac{1}{4 n \alpha}}\Big].
  \end{align*}
  This is valid for any finite $n$, but it is conservative, and the length of the confidence interval is typically larger than the asymptotic one.

  Using more advanced probability deviation inequalities, we can sometimes achieve ``best-of-both-worlds'' confidence intervals that are valid for finite $n$ and also have good asymptotic properties.
\end{itemize}


\end{document}
